% !TEX root = ../thesis.tex

\chapter{Požiadavky na systém}\label{ch:functional-requirements}

Táto kapitola formuluje počiatočné funkčné požiadavky na navrhovaný informačný systém. Požiadavky sú rozdelené do logických blokov a sú písané stručne, aby boli jednoznačné pre návrh, implementáciu aj testovanie.

\section{Funkčné požiadavky}

Funkčné požiadavky určujú, aké činnosti má systém vykonávať a aké procesy prevádzky má pokrývať. Tento blok slúži ako základ pre návrh jednotlivých modulov aplikácie a pre následné overenie, či riešenie plní očakávané funkcie.

\subsection{Evidencia zariadení}

\begin{enumerate}
\item Systém musí umožniť založiť kartu zariadenia s minimálne týmito údajmi: identifikátor, typ, umiestnenie, stav a dátum uvedenia do prevádzky.
\item Systém musí umožniť upravovať údaje zariadenia a evidovať históriu vykonaných zmien.
\item Systém musí umožniť vyhľadávanie a filtrovanie zariadení podľa typu, lokality, stavu a zodpovednej osoby.
\item Systém musí evidovať väzby medzi zariadeniami (napr. zariadenie a jeho komponenty).
\end{enumerate}

\subsection{Revízie}

\begin{enumerate}
\item Systém musí umožniť naplánovať revíziu zariadenia s termínom, rozsahom a pridelenou zodpovednou osobou.
\item Systém musí umožniť zaznamenať výsledok revízie vrátane stavu „vyhovel/nevyhovel“, poznámky a príloh.
\item Systém musí umožniť viesť úplnú históriu revízií pre každé zariadenie v časovom poradí.
\item Systém musí označiť revízie po termíne a zobraziť ich v samostatnom prehľade.
\end{enumerate}

\subsection{Servis}

\begin{enumerate}
\item Systém musí umožniť založiť servisný zásah pre zariadenie na základe poruchy, revízie alebo manuálne.
\item Systém musí evidovať priebeh zásahu: dátum začiatku, dátum ukončenia, vykonané práce, použitý materiál a zodpovednú osobu.
\item Systém musí umožniť meniť stav zásahu (nový, v riešení, ukončený, zrušený).
\item Systém musí prepojiť servisný zásah s konkrétnou revíziou alebo incidentom, ak taká väzba existuje.
\end{enumerate}

\subsection{Notifikácie}

\begin{enumerate}
\item Systém musí automaticky odoslať notifikáciu pred plánovaným termínom revízie.
\item Systém musí automaticky odoslať notifikáciu pri prekročení termínu revízie alebo servisného zásahu.
\item Systém musí umožniť nastaviť minimálne dva kanály notifikácií: e-mail a interné oznámenie v systéme.
\item Systém musí umožniť konfiguráciu príjemcov notifikácií podľa roly alebo zodpovednosti za zariadenie.
\end{enumerate}

\subsection{Reporty}

\begin{enumerate}
\item Systém musí umožniť vygenerovať report revízií za zvolené obdobie.
\item Systém musí umožniť vygenerovať report servisných zásahov podľa typu zariadenia, lokality a stavu.
\item Systém musí umožniť export reportov minimálne do formátov PDF a CSV.
\item Systém musí umožniť uložiť report do archívu a spätne ho dohľadať podľa dátumu a typu reportu.
\end{enumerate}

\section{Nefunkčné požiadavky}

Nefunkčné požiadavky určujú podmienky, za akých majú byť funkcie systému poskytované, a zabezpečujú kvalitu, spoľahlivosť a bezpečnosť riešenia.

\subsection{Bezpečnosť (Security)}

Systém musí garantovať ochranu dát a bezpečný prístup k informáciám:

\begin{itemize}
    \item \textbf{Autentifikácia a autorizácia} -- Systém musí vyžadovať autentifikáciu používateľa a autorizáciu prístupu podľa pridelenej roly (administrátor, správca siete, technik, revízny pracovník, manažér).
    \item \textbf{Šifrovaný prenos dát} -- Systém musí zabezpečiť šifrovaný prenos dát medzi klientom a serverom pomocou protokolu HTTPS.
    \item \textbf{Auditný záznam} -- Systém musí zaznamenávať dôležité operácie používateľov do auditného záznamu tak, aby bolo možné spätne overiť vykonané zmeny.
    \item \textbf{Zabezpečená správa prístupových práv} -- Systém musí umožniť centrálnu správu používateľov, ich rolí a prístupových práv.
\end{itemize}

\subsection{Dostupnosť a výkon}

Systém musí byť používateľsky prístupný, spoľahlivý a výkonný:

\begin{itemize}
    \item \textbf{Responzívne webové rozhranie} -- Systém musí poskytovať responzívne webové rozhranie použiteľné minimálne na desktopových aj mobilných zariadeniach.
    \item \textbf{Stabilná prevádzka} -- Systém musí poskytovať stabilnú prevádzku pri bežnej súbežnej práci viacerých používateľov bez straty dát.
    \item \textbf{Pravidelné zálohovanie dát} -- Systém musí byť navrhnutý tak, aby podporoval pravidelné zálohovanie dát.
    \item \textbf{Obnova po zlyhaní} -- Systém musí umožniť obnovu po zlyhaní bez straty kritických informácií.
\end{itemize}

\subsection{Archivácia a export dát}

Systém musí zabezpečiť dlhodobé uchovávanie údajov a kompatibilitu s ďalšími nástrojmi:

\begin{itemize}
    \item \textbf{Export reportov} -- Systém musí umožniť export reportov minimálne do formátov PDF a CSV.
    \item \textbf{Uchovávanie histórie} -- Systém musí uchovávať históriu zmien zariadení, revízií a servisných zásahov.
    \item \textbf{Archivácia protokolov} -- Systém musí umožniť uložiť report do archívu a spätne ho dohľadať podľa dátumu a typu reportu.
    \item \textbf{Kompatibilita} -- Systém musí umožniť export dát v štandardných formátoch tak, aby boli použiteľné v ďalších nástrojoch organizácie (napr. GIS, SCADA, CMMS).
\end{itemize}

\subsection{Rozšíriteľnosť a údržba}

Systém musí byť flexibilný a pripravený na budúce rozšírenia:

\begin{itemize}
    \item \textbf{Modulárny návrh} -- Systém musí byť navrhnutý modulárne, aby bolo možné jednoducho rozširovať funkcionalitu o nové typy zariadení, reportov a notifikácií.
    \item \textbf{Podpora nových typov zariadení} -- Systém musí umožniť pridávanie nových typov filtračných a distribučných zariadení bez nutnosti zásadných zmien architektúry.
    \item \textbf{Flexibilná konfigurácia notifikácií a reportov} -- Systém musí umožniť nastavenie vlastných pravidiel pre notifikácie a generovanie špecifických reportov.
    \item \textbf{Integrácia s existujúcimi systémami} -- Systém musí umožniť integráciu s existujúcimi systémami prevádzkovanými vo vodárenskej organizácii (napr. SCADA, GIS databázy, CMMS/EAM).
\end{itemize}


\section{Používatelia systému (Aktéri)}

Systém je určený pre viacero skupín používateľov s odlišnými úlohami a oprávneniami. Každá rola zodpovedá špecifickým potrebám pri správe distribučnej infraštruktúry a evidencii revízií.

\subsection{Administrátor}

Administrátor je zodpovedný za správu celého systému a jeho používateľov:

\begin{itemize}
    \item Spravuje používateľské účty a priradenie rolí.
    \item Nastavuje prístupové práva jednotlivým používateľom alebo skupinám.
    \item Konfiguruje systémové parametre (notifikácie, integrácie, reporty).
    \item Monitoruje auditné záznamy a zabezpečuje súlad s internými bezpečnostnými politikami.
\end{itemize}

\subsection{Správca siete}

Správca siete zabezpečuje plynulú prevádzku distribučnej infraštruktúry:

\begin{itemize}
    \item Eviduje zariadenia (pumpy, filtre, ventily, meracie zariadenia).
    \item Plánuje revízie a priradí zodpovedné osoby.
    \item Monitoruje stav zariadení a sleduje históriu zmien.
    \item Vytvára prevádzkové reporty a vyhodnocuje technický stav infraštruktúry.
\end{itemize}

\subsection{Technik}

Technik vykonáva servisné zásahy a údržbu zariadení:

\begin{itemize}
    \item Zaznamenáva servisné zásahy, opravy a výmeny komponentov.
    \item Dokumentuje použitý materiál a čas potrebný na vykonanie práce.
    \item Aktualizuje stav zariadenia po vykonaní údržby.
    \item Reaguje na notifikácie o poruchách alebo neplánovaných udalostiach.
\end{itemize}

\subsection{Revízny pracovník}

Revízny pracovník vykonáva pravidelné revízie filtračných zariadení:

\begin{itemize}
    \item Vykonáva revízie podľa harmonogramu a legislatívnych požiadaviek.
    \item Vypĺňa revízne protokoly a zaznamenáva zistené nedostatky.
    \item Navrhuje nápravné opatrenia pri identifikácii problémov.
    \item Prikladá fotodokumentáciu a ďalšie prílohy k revíznym záznamom.
\end{itemize}

\subsection{Manažér}

Manažér má strategický prehľad o stave systému a zabezpečuje kontrolu procesov:

\begin{itemize}
    \item Získava prehľad o celkovom stave zariadení a plnení revíznych termínov.
    \item Generuje a analyzuje súhrnné reporty pre rozhodovanie a audit.
    \item Kontroluje dodržiavanie harmonogramov a termínov revízií.
    \item Sleduje trendy poruchovosti, náklady na údržbu a efektivitu nápravných opatrení.
\end{itemize}

% \subsection{Pracovný tok používateľov}

% Systém podporuje komplexný pracovný tok, ktorý integruje činnosť všetkých aktérov:

% \begin{enumerate}
%     \item \textbf{Plánovanie} -- Správca siete naplánuje revízie, priradí zodpovednú osobu a nastaví termíny.
%     \item \textbf{Notifikácia} -- Systém automaticky odošle notifikáciu pred plánovaným termínom revízie.
%     \item \textbf{Vykonanie revízie} -- Revízny pracovník vykoná revíziu a zaznamená výsledok do systému (vyhovel/nevyhovel, poznámky, fotodokumentácia).
%     \item \textbf{Servisný zásah} -- Pri zistení problému systém vytvorí servisný zásah a notifikuje technika, ktorý vykoná opravu a zdokumentuje práce.
%     \item \textbf{Kontrola a reportovanie} -- Manažér pravidelně kontroluje reporty, overuje dodržiavanie termínov a vyhodnocuje trendy.
% \end{enumerate}